% ==========================================
% 二维均匀分布与几何概型
% ==========================================

\vspace{1.5em}
\subsubsection*{1. 二维均匀分布的几何概型解法}

\begin{例题}[雷达屏幕上的均匀分布]
	雷达的圆形屏幕半径为 $R$，设目标出现点 $(X,Y)$ 在屏幕上服从均匀分布。
	
	(1) 求 $P\{Y>0 \mid Y>X\}$；
	
	(2) 设 $M = \max\{X,Y\}$，求 $P\{M>0\}$。
\end{例题}

\textbf{解答}\quad

因为点 $(X,Y)$ 在半径为 $R$ 的圆 $D: x^2 + y^2 \le R^2$ 上服从均匀分布，所以这是一个典型的\textbf{几何概型}问题。对于圆内的任意区域 $G$，点落在 $G$ 内的概率等于 $G$ 的面积与圆面积的比值，即 $P((X,Y) \in G) = \frac{S_G}{\pi R^2}$。

\textbf{(1) 求 $P\{Y>0 \mid Y>X\}$} \\
根据条件概率公式：
$$ P\{Y>0 \mid Y>X\} = \frac{P\{Y>0, Y>X\}}{P\{Y>X\}} $$
首先计算分母 $P\{Y>X\}$：\\
区域 $Y>X$ 是直线 $y=x$ 上方的半圆。根据对称性，其面积正好是整个圆面积的一半，因此：
$$ P\{Y>X\} = \frac{1}{2} $$
接着计算分子 $P\{Y>0, Y>X\}$：\\
该事件要求点同时满足 $y>0$（在 $x$ 轴上方）和 $y>x$（在直线 $y=x$ 上方）。\\
在极坐标系下，$y>0$ 对应角度 $\theta \in (0, \pi)$，$y>x$ 对应角度 $\theta \in (\frac{\pi}{4}, \frac{5\pi}{4})$。\\
两者的交集为 $\theta \in (\frac{\pi}{4}, \pi)$。这个扇形的圆心角为 $\pi - \frac{\pi}{4} = \frac{3\pi}{4}$，占整个圆周角 $2\pi$ 的 $\frac{3}{8}$。因此：
$$ P\{Y>0, Y>X\} = \frac{3}{8} $$
代入条件概率公式：
$$ P\{Y>0 \mid Y>X\} = \frac{3/8}{1/2} = \frac{3}{4} $$

\vspace{0.5em}
\textbf{(2) 求 $P\{M>0\}$} \\
已知 $M = \max\{X,Y\}$。直接求 $M>0$ 比较复杂，我们可以利用\textbf{对立事件}来转化：
$$ P\{M>0\} = 1 - P\{M \le 0\} $$
而最大值小于等于 0，意味着两个变量\textbf{同时}小于等于 0：
$$ P\{M \le 0\} = P\{\max\{X,Y\} \le 0\} = P\{X \le 0, Y \le 0\} $$
区域 $X \le 0, Y \le 0$ 正好是圆在\textbf{第三象限}的部分。显然，它的面积是整个圆面积的 $\frac{1}{4}$。因此：
$$ P\{X \le 0, Y \le 0\} = \frac{1}{4} $$
所以：
$$ P\{M>0\} = 1 - \frac{1}{4} = \frac{3}{4} $$